% Capítulo: La soberanía de Dios (generado desde ../06_Soberano/06_Soberano_capitulo.tex)
% Texto ancla: Daniel 4:34–35. Clasificación: Identidad; incomunicable como gobierno de todas las cosas, comunicable como reflejo (señorío delegado, Sal 8:6).

\section{Un atributo de identidad}

La definición de partida de este estudio cabe en una línea: \destacar{Dios gobierna absolutamente todas las cosas según su voluntad}. Cada una de sus partes tiene un texto detrás. ``Gobierna'': ``el Altísimo gobierna el reino de los hombres'' (Dn 4:17). ``Todas las cosas'': ``su reino domina sobre todos'' (Sal 103:19). ``Según su voluntad'': Dios ``hace todas las cosas según el designio de su voluntad'' (Ef 1:11). El salmista lo dice en forma de contraste con los ídolos:

\hebverso{וֵאלֹהֵינוּ בַשָּׁמָיִם כֹּל אֲשֶׁר־חָפֵץ עָשָׂה}
\begin{verso}{Salmo 115:3}
Nuestro Dios está en los cielos; Todo lo que quiso ha hecho.
\end{verso}

\begin{filologia}{Sal 115:3}
\heb{בַשָּׁמָיִם} \emph{vashamáyim} & \heb{שׁמים} & ``en los cielos'': el lugar del trono (Sal 103:19). \\
\heb{חָפֵץ} \emph{jafets} & \heb{חפץ} & ``quiso, se complació en''. Es el deseo libre de quien no depende de nada. El mismo verbo vuelve en Sal 135:6, Is 46:10 y Jer 9:24 (``estas cosas quiero''). \\
\heb{עָשָׂה} \emph{asá} & \heb{עשׂה} & ``hizo'': lo que Dios quiere, lo hace. Los ídolos, en cambio, ``tienen boca, mas no hablan'' (115:5). \\
\end{filologia}

La soberanía responde a la pregunta \destacar{¿quién es Dios?}: es el que reina. No es una función que Dios cumpla a veces, sino lo que él es; por eso esta serie la cuenta entre los atributos de identidad. Se apoya en otros: si Dios es soberano, lo es porque es eterno (``cuyo dominio es sempiterno'', Dn 4:34), omnipotente y sabio. Y alcanza tres ámbitos que la Escritura nombra por separado: la naturaleza (``todo lo que \SENOR{} quiere, lo hace, en los cielos y en la tierra, en los mares y en todos los abismos'', Sal 135:6), las naciones (``quita reyes, y pone reyes'', Dn 2:21) y los corazones (``como los repartimientos de las aguas, así está el corazón del rey en la mano de \SENOR'', Pr 21:1).

\precision{clasificación propuesta: atributo de Identidad; incomunicable como gobierno de todas las cosas, comunicable como reflejo en el señorío que Dios delegó al ser humano: ``Le hiciste señorear sobre las obras de tus manos'' (Sal 8:6 [8:7]; cf. Gn 1:28). Queda sujeta a la revisión del autor de la serie.}

\subsection{La pregunta del dolor}
Si Dios gobierna todo, ¿por qué hay tanto dolor? Y si una gran parte del dolor la causa el mismo ser humano, ¿sigue Dios siendo soberano? C. S. Lewis planteó el problema en su forma clásica en \emph{El problema del dolor} (1940), en el capítulo titulado ``La omnipotencia divina'': si Dios fuera bueno, desearía que sus criaturas fueran perfectamente felices, y si fuera todopoderoso, sería capaz de hacer lo que desea; pero las criaturas no son felices; luego Dios carecería de bondad, de poder o de ambos. Este capítulo toma esa pregunta en serio y la lleva a los textos: a Daniel, a Isaías, a José, a los sabios de Israel y a Pablo.

\subsection{“Soberano” en la Biblia}
La palabra ``soberano'' es escasa en la RVR1960, pero la idea está en todas partes.

\begin{center}\small\renewcommand{\arraystretch}{1.18}
\begin{tabular}{@{}L{2.6cm}L{4.4cm}L{6.6cm}@{}}
\toprule
\textbf{Texto} & \textbf{RVR1960} & \textbf{Original} \\ \midrule
Job 31:28 & ``el Dios soberano'' & \heb{אֵל מִמָּעַל} \emph{El mimmá'al}, ``Dios de arriba'' \\
Ez 38:2, 3; 39:1 & ``príncipe soberano'' (Gog) & \heb{נְשִׂיא רֹאשׁ} \emph{nesí rosh} \\
Hch 4:24 & ``Soberano Señor'' & \grc{Δέσποτα} \\
1 Ti 6:15 & ``el bienaventurado y solo Soberano'' & \grc{ὁ μακάριος καὶ μόνος δυνάστης} \\
Jud 4 & ``Dios el único soberano'' & \grc{τὸν μόνον δεσπότην} \\
Ap 1:5 & ``el soberano de los reyes de la tierra'' & \grc{ὁ ἄρχων τῶν βασιλέων τῆς γῆς} \\
\bottomrule
\end{tabular}\par
{\footnotesize\color{gris}\textit{Tabla 1.} ``Soberano'' en la RVR1960. \textit{Nota.} Conteo propio, sin títulos de sección: 8 apariciones en 8 versículos; ``soberanía'' solo figura en el título de sección de Job 26.}
\end{center}

\precision{en Jud 4 el griego dice ``negando a nuestro único Soberano y Señor Jesucristo''; la RVR1960, siguiendo el texto bizantino (que añade \grc{θεόν}), lee ``niegan a Dios el único soberano, y a nuestro Señor Jesucristo''. En Ap 1:5 el título se aplica a Jesucristo.}

El vocabulario de la soberanía está en los nombres de Dios (\emph{Adonay}, \emph{Elyón}), en sus verbos (reinar, gobernar, dominar, querer) y en fórmulas como ``\SENOR{} reina'' o ``todo lo que quiso ha hecho''. Se estudia en la sección sobre el Nombre.

%===================================================================
\section{Texto ancla: Daniel 4}

\subsection{Un rey que tuvo que aprender}
Daniel 4 es una carta de Nabucodonosor ``a todos los pueblos, naciones y lenguas que moran en toda la tierra'' (4:1), escrita en primera persona por el rey más poderoso de su tiempo. Su tema es uno: \destacar{quién gobierna de verdad}. El rey sueña con un árbol que llega hasta el cielo y es cortado (4:10--17); Daniel le advierte; un año después el rey se pasea por su palacio y se jacta:

\hebverso{הֲלָא דָא־הִיא בָּבֶל רַבְּתָא דִּי־אֲנָה בֱנַיְתַהּ}
\begin{verso}{Daniel 4:30 [4:27]}
habló el rey y dijo: ¿No es ésta la gran Babilonia que yo edifiqué para casa real con la fuerza de mi poder, y para gloria de mi majestad?
\end{verso}

``Aún estaba la palabra en la boca del rey'' cuando cae la sentencia (4:31). Nabucodonosor pierde la razón, vive como un animal ``hasta que su pelo creció como plumas de águila'' (4:33) y, al final, levanta los ojos al cielo.

\precision{dn 2:4b--7:28 está escrito en arameo. La numeración española de Dn 4 va tres versículos adelante de la aramea (4:1--3 en español corresponde a 3:31--33 en arameo).}

\subsection{El estribillo: “hasta que conozcas”}
El capítulo repite cuatro veces la misma lección, como un estribillo:

\begin{center}\small\renewcommand{\arraystretch}{1.18}
\begin{tabular}{@{}L{2.0cm}L{7.4cm}L{5.0cm}@{}}
\toprule
\textbf{Texto} & \textbf{RVR1960} & \textbf{Arameo} \\ \midrule
4:17 [4:14] & ``para que conozcan los vivientes que el Altísimo gobierna el reino de los hombres, y que a quien él quiere lo da'' & \heb{דִּי־שַׁלִּיט עִלָּיָא בְּמַלְכוּת אֱנָשָׁא} \\
4:25 [4:22] & ``hasta que conozcas que el Altísimo tiene dominio en el reino de los hombres, y que lo da a quien él quiere'' & \heb{עַד דִּי־תִנְדַּע דִּי־שַׁלִּיט עִלָּיָא} \\
4:26 [4:23] & ``luego que reconozcas que el cielo gobierna'' & \heb{דִּי שַׁלִּטִן שְׁמַיָּא} \\
4:32 [4:29] & ``hasta que reconozcas que el Altísimo tiene el dominio en el reino de los hombres, y lo da a quien él quiere'' & \heb{עַד דִּי־תִנְדַּע דִּי־שַׁלִּיט עִלָּיָא} \\
\bottomrule
\end{tabular}\par
{\footnotesize\color{gris}\textit{Tabla 2.} El estribillo de Daniel 4. \textit{Nota.} Elaboración propia (RVR1960 y BHS).}
\end{center}

La forma literaria es parte del mensaje: el capítulo está construido para que el lector aprenda lo mismo que el rey. El verbo es \heb{שַׁלִּיט} \emph{shallit}, ``es soberano, tiene dominio''; el sujeto, \heb{עִלָּיָא} \emph{Illayá}, ``el Altísimo''.

\precision{en 4:26 ``el cielo'' sustituye a ``Dios'': ``el cielo gobierna''. El judaísmo posterior generalizó ese uso (\emph{malkhut shamayim}, ``reino de los cielos''), que Mateo conserva en su evangelio.}

\subsection{La confesión}
Al final de los ``siete tiempos'', el rey confiesa:

\hebverso{וּלְעִלָּיָא בָּרְכֵת וּלְחַי עָלְמָא שַׁבְּחֵת וְהַדְּרֵת \\ דִּי שָׁלְטָנֵהּ שָׁלְטָן עָלַם וּמַלְכוּתֵהּ עִם־דָּר וְדָר \\ וּכְמִצְבְּיֵהּ עָבֵד בְּחֵיל שְׁמַיָּא וְדָאֲרֵי אַרְעָא \\ וְלָא אִיתַי דִּי־יְמַחֵא בִידֵהּ וְיֵאמַר לֵהּ מָה עֲבַדְתְּ}
\begin{verso}{Daniel 4:34–35 [4:31–32]}
Mas al fin del tiempo yo Nabucodonosor alcé mis ojos al cielo, y mi razón me fue devuelta; y bendije al Altísimo, y alabé y glorifiqué al que vive para siempre, cuyo dominio es sempiterno, y su reino por todas las edades. Todos los habitantes de la tierra son considerados como nada; y él hace según su voluntad en el ejército del cielo, y en los habitantes de la tierra, y no hay quien detenga su mano, y le diga: ¿Qué haces?
\end{verso}

\begin{filologia}{Dn 4:34--35}
\heb{שָׁלְטָנֵהּ שָׁלְטָן עָלַם} & \heb{שׁלט} & ``su dominio (es) dominio eterno'': \emph{sholtán}, ``dominio, imperio''. Es la misma raíz de \emph{shallit} (4:17); el rey que ``dominaba'' confiesa un dominio que no termina. \\
\heb{כְּלָה חֲשִׁיבִין} \emph{kelá jashivín} & \heb{חשׁב} & ``como nada son considerados'': los habitantes de la tierra no ``cuentan'' frente a él. \\
\heb{וּכְמִצְבְּיֵהּ עָבֵד} & \heb{צבא} & ``y conforme a su voluntad hace'': \emph{tsevá}, ``querer'' (Strong H6634), el mismo verbo de ``a quien él quiere lo da'' (4:17, 25, 32). \\
\heb{בְּחֵיל שְׁמַיָּא} & \heb{חיל} & ``en el ejército del cielo'': los poderes celestiales. \\
\heb{דִּי־יְמַחֵא בִידֵהּ} & \heb{מחא} & ``que golpee su mano'': gesto de quien detiene o reprocha. Nadie puede frenar la mano de Dios. \\
\heb{מָה עֲבַדְתְּ} & \heb{עבד} & ``¿qué has hecho?'': nadie le pide cuentas; es la fórmula de Job 9:12, Is 45:9 y Ec 8:4 (allí, del rey humano). \\
\end{filologia}

La versión griega de Teodoción traduce ``conforme a su voluntad hace'' con \grc{κατὰ τὸ θέλημα αὐτοῦ ποιεῖ}, y ``Altísimo'' con \grc{ὕψιστος}: el vocabulario que reaparecerá en Pablo (Ef 1:11).

\hebverso{לְמֶלֶךְ שְׁמַיָּא דִּי כָל־מַעֲבָדוֹהִי קְשֹׁט וְאֹרְחָתֵהּ דִּין}
\begin{verso}{Daniel 4:37 [4:34]}
Ahora yo Nabucodonosor alabo, engrandezco y glorifico al Rey del cielo, porque todas sus obras son verdaderas, y sus caminos justos; y él puede humillar a los que andan con soberbia.
\end{verso}

\begin{filologia}{Dn 4:37}
\heb{מֶלֶךְ שְׁמַיָּא} \emph{mélej shemayá} & \heb{מלך} & ``Rey del cielo'': el que se creía rey reconoce a otro Rey. \\
\heb{קְשֹׁט} \emph{qeshot} & \heb{קשׁט} & ``verdad'': sus obras corresponden a lo que él es. \\
\heb{דִּין} \emph{din} & \heb{דין} & ``juicio, justicia'': sus caminos son rectos. \\
\end{filologia}

La confesión no deja la soberanía sola: el Rey del cielo actúa con \destacar{verdad} y con \destacar{justicia}. Y el medio por el que el rey aprendió fue el dolor de la humillación: ``siete tiempos pasarán sobre ti, hasta que conozcas'' (4:25).

\precision{que la humillación de Nabucodonosor fue un instrumento pedagógico de la soberanía lo dice el propio texto (``hasta que conozcas''). Su relación con el ``megáfono'' de C. S. Lewis (véase la sección sobre Lewis) es una conexión propia de esta serie.}

%===================================================================
\section{El Nombre que revela la soberanía}

En la Biblia el nombre revela quién es Dios (Éx 3:15; 34:5--7). Cada nombre destaca un aspecto de su ser. El nombre que más directamente expresa la soberanía es \destacar{\emph{Adonay}}: Dios es el Señor, el Dueño que dispone de lo suyo. \emph{Elyón}, ``el Altísimo'', añade la altura: Dios está por encima de todo poder. A ellos se suma la fórmula de entronización ``el \SENOR{} reina''.

\subsection{\emph{Adonay}: el Dueño}
El sustantivo \heb{אָדוֹן} \emph{adón} significa ``señor, amo, dueño''; Strong lo define como ``soberano, controlador (humano o divino)''. Se dice de un amo con sus siervos y de un rey con sus súbditos. La forma \heb{אֲדֹנָי} \emph{Adonay}, con \emph{qamets} final, es enfática y se reserva a Dios: ``el Señor''. Malaquías muestra lo que implica:

\hebverso{וְאִם־אֲדוֹנִים אָנִי אַיֵּה מוֹרָאִי}
\begin{verso}{Malaquías 1:6}
El hijo honra al padre, y el siervo a su señor. Si, pues, soy yo padre, ¿dónde está mi honra? y si soy señor, ¿dónde está mi temor? dice \SENOR{} de los ejércitos a vosotros, oh sacerdotes, que menospreciáis mi nombre.
\end{verso}

\begin{filologia}{Mal 1:6}
\heb{אֲדֹנָיו} \emph{adonav} & \heb{אדן} & ``su señor'': el amo del siervo. \\
\heb{אֲדוֹנִים} \emph{adoním} & \heb{אדן} & ``señor'', en plural de majestad aplicado a Dios. \\
\heb{מוֹרָאִי} \emph{moraí} & \heb{ירא} & ``mi temor'': al Señor le corresponde reverencia. \\
\end{filologia}

\begin{center}\small\renewcommand{\arraystretch}{1.18}
\begin{tabular}{@{}L{3.4cm}L{10.2cm}@{}}
\toprule
\textbf{Forma} & \textbf{Uso} \\ \midrule
\heb{אֲדֹנִי} \emph{adoní} & ``mi señor'', dicho a un hombre: unas 200 veces. \\
\heb{אֲדֹנָי} \emph{Adonay} & ``el Señor'', de Dios: 437 veces en 422 versículos. \\
\heb{אֲדֹנָי יֱהוִה} & ``Señor \SENOR'' o ``\SENOR{} el Señor'' en la RVR1960: 301 veces, 210 en Ezequiel. \\
\heb{אֲדֹנֵינוּ} \emph{adonenu} & ``Señor nuestro'' (Sal 8:1 [8:2]). \\
\bottomrule
\end{tabular}\par
{\footnotesize\color{gris}\textit{Tabla 3.} Las formas de \emph{adón}. \textit{Nota.} Conteo propio en la BHS por forma consonántica y vocal final; verificar en e-Sword.}
\end{center}

El contraste entre señorío delegado y señorío propio se ve en José. Faraón lo pone ``por señor de toda su casa'' (\heb{וּלְאָדוֹן לְכָל־בֵּיתוֹ}, Gn 45:8): es \emph{adón}, pero porque otro lo puso. El señorío de Dios no lo recibe de nadie.

\precision{al leer en voz alta, el judaísmo pronuncia \emph{Adonay} donde está escrito el Nombre \heb{יהוה}. Por eso esta serie escribe \SENOR{} en lugar de ``Jehová'': en la lectura de Israel, el Nombre propio de Dios se dice con el título de su soberanía.}

\subsection{“Señor \SENOR”: Abram y el Dueño}
La primera vez que la Biblia une \emph{Adonay} y el Nombre es en la noche de la espera de Abram:

\hebverso{וַיֹּאמֶר אַבְרָם אֲדֹנָי יֱהוִה מַה־תִּתֶּן־לִי וְאָנֹכִי הוֹלֵךְ עֲרִירִי}
\begin{verso}{Génesis 15:2}
Y respondió Abram: Señor \SENOR, ¿qué me darás, siendo así que ando sin hijo, y el mayordomo de mi casa es ese damasceno Eliezer?
\end{verso}

\begin{filologia}{Gn 15:2}
\heb{אֲדֹנָי יֱהוִה} & \heb{אדן} & ``Señor \SENOR'' (RVR1960: ``Señor Jehová''). Los masoretas pusieron en el Nombre las vocales de \emph{Elohim}, para que se lea \emph{Adonay Elohim} y no se repita ``Adonay''. \\
\heb{עֲרִירִי} \emph{aririí} & \heb{ערר} & ``sin hijos'', ``desnudo'' de descendencia. \\
\textbf{LXX} & --- & \grc{Δέσποτα, τί μοι δώσεις;} ``Soberano, ¿qué me darás?''; en 15:8, \grc{Δέσποτα κύριε}. \\
\end{filologia}

Abram no tiene hijo ni heredero, y la promesa parece imposible. Justamente ahí llama a Dios \destacar{Dueño}. La Septuaginta tradujo ese primer \emph{Adonay YHWH} con \grc{δέσποτα}, ``amo con autoridad plena'', y siglos después la iglesia de Jerusalén, amenazada, ora con la misma palabra: ``Soberano Señor, tú eres el Dios que hiciste el cielo y la tierra'' (Hch 4:24).

\subsection{\emph{Elyón}: el Altísimo}
El segundo nombre de la soberanía aparece en el encuentro de Abram con Melquisedec, después de la victoria sobre los cuatro reyes de oriente:

\hebverso{בָּרוּךְ אַבְרָם לְאֵל עֶלְיוֹן קֹנֵה שָׁמַיִם וָאָרֶץ \\ וּבָרוּךְ אֵל עֶלְיוֹן אֲשֶׁר־מִגֵּן צָרֶיךָ בְּיָדֶךָ}
\begin{verso}{Génesis 14:19–20}
y le bendijo, diciendo: Bendito sea Abram del Dios Altísimo, creador de los cielos y de la tierra; y bendito sea el Dios Altísimo, que entregó tus enemigos en tu mano. Y le dio Abram los diezmos de todo.
\end{verso}

\begin{filologia}{Gn 14:19--20}
\heb{אֵל עֶלְיוֹן} \emph{El Elyón} & \heb{עלה} & ``Dios Altísimo'': de \emph{alá}, ``subir''; el que está más alto. \\
\heb{קֹנֵה} \emph{qoné} & \heb{קנה} & ``creador'' (RVR1960); también ``poseedor, el que adquiere''. Los léxicos registran ambos sentidos (BDB; HALOT): Dios es dueño de lo que hizo. \\
\heb{מִגֵּן} \emph{miggén} & \heb{מגן} & ``entregó'': Dios es quien dio la victoria, no la espada de Abram. \\
\end{filologia}

\emph{Elyón} es el Dios que reparte el mapa del mundo: ``Cuando el Altísimo hizo heredar a las naciones, cuando hizo dividir a los hijos de los hombres, estableció los límites de los pueblos'' (Dt 32:8). Los salmos lo unen a la realeza: ``Porque \SENOR{} el Altísimo es temible; Rey grande sobre toda la tierra'' (Sal 47:2 [47:3]). En Daniel, en arameo, \heb{עִלָּיָא} \emph{Illayá} es el título dominante del capítulo 4.

\subsection{“El \SENOR{} reina”}
A los nombres se suma una fórmula: \heb{יְהוָה מָלָךְ} \emph{Adonay malaj}, ``el \SENOR{} reina'' (Sal 93:1; 96:10; 97:1; 99:1; 1 Cr 16:31), que ya resuena en el cántico del mar: ``\SENOR{} reinará eternamente y para siempre'' (Éx 15:18). Y otra, de gobierno: ``\SENOR{} estableció en los cielos su trono, y su reino domina sobre todos'' (Sal 103:19), donde ``domina'' es \heb{מָשָׁלָה} \emph{mashalá}. La Septuaginta traduce ese verbo con \grc{δεσπόζει}, de la familia de \grc{δεσπότης}.

\subsection{El espectro completo}
\begin{center}\small\renewcommand{\arraystretch}{1.15}
\begin{tabular}{@{}L{3.6cm}L{3.4cm}L{2.6cm}L{4.2cm}@{}}
\toprule
\textbf{Término} & \textbf{Sentido} & \textbf{RVR1960} & \textbf{Textos} \\ \midrule
\heb{אֲדֹנָי} \emph{Adonay} & el Señor, el Dueño & Señor & Gn 15:2; Sal 8:1; Is 6:1; Lm 3:37 \\
\heb{עֶלְיוֹן} \emph{Elyón} & el Altísimo & Altísimo & Gn 14:18--22; Dt 32:8; Sal 47:2; 91:1 \\
\heb{מָלַךְ} \emph{malaj} & reinar & reinar & Éx 15:18; Sal 93:1; 97:1 \\
\heb{מָשַׁל} \emph{mashal} & gobernar, dominar & dominar, señorear & 1 Cr 29:12; Sal 103:19; Sal 8:6 \\
\heb{חָפֵץ} \emph{jafets} & querer, complacerse & querer & Sal 115:3; 135:6; Is 46:10 \\
\heb{עֵצָה} \emph{etsá} & consejo, plan & consejo & Is 46:10; Pr 19:21 \\ \midrule
\heb{שַׁלִּיט} \emph{shallit} (arameo) & soberano, que domina & gobierna, tiene dominio & Dn 4:17, 25, 26, 32 \\
\heb{שָׁלְטָן} \emph{sholtán} (arameo) & dominio, imperio & dominio, señorío & Dn 4:3, 34; 7:14 \\
\heb{עִלָּיָא} \emph{Illayá} (arameo) & el Altísimo & Altísimo & Dn 4:2, 17, 34; 5:18, 21 \\ \midrule
\grc{δεσπότης} (10) & amo con autoridad plena & Señor, Soberano, amo & Lc 2:29; Hch 4:24; Jud 4; Ap 6:10 \\
\grc{δυνάστης} (3) & potentado & Soberano, poderoso & 1 Ti 6:15; Lc 1:52 \\
\grc{κύριος} & señor; en la LXX, el Nombre & Señor & Flp 2:11; Ro 10:9 \\
\grc{ὕψιστος} (13) & altísimo & Altísimo & Lc 1:32; Hch 7:48 \\
\grc{βουλή} (12) & consejo, designio & designio, consejo & Hch 2:23; 4:28; Ef 1:11 \\
\grc{θέλημα} (62) & voluntad & voluntad & Ef 1:5, 9, 11; Ro 9:19 \\
\grc{πρόθεσις} (12) & propósito & propósito & Ro 8:28; 9:11; Ef 1:11 \\
\grc{ἐνεργέω} (21) & obrar eficazmente & hacer, operar & Ef 1:11, 20; Flp 2:13 \\
\bottomrule
\end{tabular}\par
{\footnotesize\color{gris}\textit{Tabla 4.} El vocabulario de la soberanía. \textit{Nota.} Conteo propio con la BHS y NA27 con números Strong.}
\end{center}

\precision{\emph{Elyón} aparece 53 veces en la BHS; 31 son título divino y las demás significan ``superior, alto'' en sentido físico (Gn 40:17; Jos 16:5). \emph{Illayá} aparece 10 veces, y \emph{elyonín} (``santos del Altísimo'', Dn 7:18--27), 4. En la RVR1960, ``Altísimo'' aparece 57 veces en 56 versículos. De los 10 usos de \grc{δεσπότης} en el NT, 6 se refieren a Dios o a Cristo y 4 a amos humanos (1 Ti 6:1--2; Tit 2:9; 1 P 2:18). Conviene verificar los conteos en e-Sword.}

\subsection{Del hebreo al griego de Pablo}
La cadena es sencilla y decisiva. En hebreo, \heb{אֲדֹנָי} es el Dueño y la palabra con que Israel pronuncia el Nombre. La Septuaginta traduce el Nombre y \emph{Adonay} con \grc{κύριος}, y a veces con \grc{δέσποτα} (Gn 15:2). Pablo confiesa: \grc{κύριος Ἰησοῦς Χριστός}, ``Jesucristo es el Señor'' (Flp 2:11). El nombre que Israel decía en lugar del Nombre es el que Pablo confiesa de Jesús. El salmo más citado del Nuevo Testamento es, además, un salmo de soberanía:

\hebverso{נְאֻם יְהוָה לַאדֹנִי שֵׁב לִימִינִי \\ עַד־אָשִׁית אֹיְבֶיךָ הֲדֹם לְרַגְלֶיךָ}
\begin{verso}{Salmo 110:1}
\SENOR{} dijo a mi Señor: Siéntate a mi diestra, Hasta que ponga a tus enemigos por estrado de tus pies.
\end{verso}

\begin{filologia}{Sal 110:1}
\heb{נְאֻם יְהוָה} \emph{neum Adonay} & \heb{נאם} & ``oráculo del \SENOR''. \\
\heb{לַאדֹנִי} \emph{laadoní} & \heb{אדן} & ``a mi señor'' (\emph{adoní}). La LXX: \grc{τῷ κυρίῳ μου}. \\
\heb{הֲדֹם לְרַגְלֶיךָ} & \heb{הדם} & ``estrado de tus pies'': sometimiento total de los enemigos. \\
\end{filologia}

Dios entroniza a su Señor hasta someter a todos sus enemigos. Pablo lo usará en 1 Co 15:25 y Ef 1:20--22.

%===================================================================
\section{Definición}

\subsection{¿Qué es la soberanía de Dios?}
De los textos estudiados resulta que la soberanía de Dios es \destacar{gobierno}, no solo poder (``gobierna el reino de los hombres''); es \destacar{universal} (cielo, tierra, mares y abismos); es \destacar{libre}, porque hace ``lo que quiere'' sin consejero (Is 40:13); es \destacar{eterna} (``dominio sempiterno''); y es \destacar{buena}: sus obras son verdad y sus caminos justicia (Dn 4:37).

\begin{tcolorbox}[colback=azul!5!white,colframe=azul,boxrule=1pt,arc=3pt,title={\small Definición},colbacktitle=azul,coltitle=white]
La soberanía de Dios es la perfección eterna por la cual Dios, como Señor (\emph{Adonay}) y Altísimo (\emph{Elyón}), gobierna absolutamente todas las cosas según su voluntad: dispone de lo que es suyo con derecho pleno, cumple su consejo por todos los medios, aun los que el hombre usa para mal, y lo encamina todo a su gloria y al bien de su pueblo, sin dejar de ser justo ni de hacer responsable a la criatura.
\end{tcolorbox}

\precision{definición de síntesis de esta serie, a partir de la definición de partida (``Dios gobierna absolutamente todas las cosas según su voluntad'') y de los textos estudiados.}

\subsection{Lo que la soberanía no es}
\begin{center}\small\renewcommand{\arraystretch}{1.2}
\begin{tabular}{@{}L{2.6cm}L{4.8cm}L{6.4cm}@{}}
\toprule
\textbf{Idea} & \textbf{Qué afirma} & \textbf{Por qué no es la soberanía bíblica} \\ \midrule
Fatalismo & un destino ciego e impersonal decide todo & Dios es una persona que quiere, ama y escucha (Sal 115:3; Hch 4:24--31) \\
Deísmo & Dios hizo el mundo y lo dejó funcionar solo & ``todas las cosas en él subsisten'' (Col 1:17); la suerte se echa, ``mas de \SENOR{} es la decisión de ella'' (Pr 16:33) \\
Dualismo & el bien y el mal son dos poderes iguales & ``Yo \SENOR{} soy el que hago todo esto'' (Is 45:7) \\
Arbitrariedad & Dios actúa por capricho & sus caminos son ``justos'' (Dn 4:37); hace ``misericordia, juicio y justicia'' (Jer 9:24) \\
Pasividad humana & si Dios decide, el hombre no responde & ``ocupaos en vuestra salvación\ldots{} porque Dios es el que en vosotros produce'' (Flp 2:12--13) \\
\bottomrule
\end{tabular}\par
{\footnotesize\color{gris}\textit{Tabla 5.} Lo que la soberanía no es. \textit{Nota.} Elaboración propia.}
\end{center}

\hebverso{לֵב אָדָם יְחַשֵּׁב דַּרְכּוֹ וַיהוָה יָכִין צַעֲדוֹ}
\begin{verso}{Proverbios 16:9}
El corazón del hombre piensa su camino; Mas \SENOR{} endereza sus pasos.
\end{verso}

\begin{filologia}{Pr 16:9}
\heb{יְחַשֵּׁב} \emph{yejashev} & \heb{חשׁב} & ``piensa, planea'': el mismo verbo de Gn 50:20. El hombre planea de verdad. \\
\heb{יָכִין} \emph{yajín} & \heb{כון} & ``establece, afirma'': Dios fija el paso. \\
\end{filologia}

%===================================================================
\section{Teología}

\subsection{Supremacía y soberanía (Pink)}
A. W. Pink dedica dos capítulos seguidos a este tema: ``La supremacía de Dios'' y ``La soberanía de Dios'' (2020, caps. 5--6). Define la soberanía como ``el ejercicio de su Supremacía'': siendo infinitamente elevado sobre la criatura más elevada, Dios ``hace como le place, sólo como le place y siempre como le place'' (Pink, 2020, p. 30). Su texto de apoyo es el mismo de este capítulo, Dn 4:35, junto con Is 46:10 (``Mi consejo permanecerá, y haré todo lo que quiero'') y Ef 1:11.

Pink cita a C. H. Spurgeon, en un sermón sobre Mateo 20:15: ``No hay atributo más reconfortante para sus hijos que el de la soberanía de Dios''. Bajo las pruebas más severas, dice Spurgeon, los creyentes creen que la soberanía ha ordenado sus aflicciones, que las gobierna y que las santificará todas (Pink, 2020, pp. 30--31). El consuelo no está en entender cada dolor, sino en saber quién está en el trono.

\precision{pink se usa en esta serie como referencia para algunos conceptos; el orden y la estructura del capítulo son propios.}

\subsection{Soberanía y responsabilidad}
Pink sostiene que la responsabilidad humana no compite con la soberanía de Dios, sino que \emph{se basa} en ella: Dios colocó a sus criaturas en el estado que quiso, les dio una ley, y por eso responden ante él (Pink, 2020, pp. 31--33). La Escritura afirma las dos cosas a la vez, muchas veces en el mismo versículo:

\begin{center}\small\renewcommand{\arraystretch}{1.18}
\begin{tabular}{@{}L{2.2cm}L{5.6cm}L{6.0cm}@{}}
\toprule
\textbf{Texto} & \textbf{Dios} & \textbf{El ser humano} \\ \midrule
Gn 50:20 & ``mas Dios lo encaminó a bien'' & ``vosotros pensasteis mal contra mí'' \\
Pr 16:9 & ``\SENOR{} endereza sus pasos'' & ``el corazón del hombre piensa su camino'' \\
Hch 2:23 & ``entregado por el determinado consejo y anticipado conocimiento de Dios'' & ``prendisteis y matasteis por manos de inicuos'' \\
Flp 2:12--13 & ``Dios es el que en vosotros produce así el querer como el hacer'' & ``ocupaos en vuestra salvación con temor y temblor'' \\
\bottomrule
\end{tabular}\par
{\footnotesize\color{gris}\textit{Tabla 6.} Soberanía y responsabilidad en el mismo texto. \textit{Nota.} Elaboración propia.}
\end{center}

La Biblia no explica \emph{cómo} se unen; muestra \emph{que} se unen. La soberanía no anula la culpa del que hace el mal ni el deber del que obedece.

\subsection{Soberano sobre el mal}
El texto más fuerte está en Isaías, frente a Ciro:

\hebverso{יוֹצֵר אוֹר וּבוֹרֵא חֹשֶׁךְ עֹשֶׂה שָׁלוֹם וּבוֹרֵא רָע \\ אֲנִי יְהוָה עֹשֶׂה כָל־אֵלֶּה}
\begin{verso}{Isaías 45:7}
que formo la luz y creo las tinieblas, que hago la paz y creo la adversidad. Yo \SENOR{} soy el que hago todo esto.
\end{verso}

\begin{filologia}{Is 45:7}
\heb{יוֹצֵר} \emph{yotser} & \heb{יצר} & ``formo'', el verbo del alfarero (cf. 45:9). \\
\heb{בּוֹרֵא} \emph{boré} & \heb{ברא} & ``creo'': el verbo de Gn 1:1 (\emph{bará}). \\
\heb{שָׁלוֹם} // \heb{רָע} & \heb{רעע} & ``paz'' // ``adversidad'': \emph{ra'} es aquí mal, desgracia, calamidad; el par es paz/calamidad, no virtud/pecado. \\
\end{filologia}

\precision{muchos intérpretes leen Is 45:7 también como respuesta al dualismo persa, que oponía un dios de la luz a uno de las tinieblas. El texto no dice que Dios sea autor del pecado: ``Dios no puede ser tentado por el mal, ni él tienta a nadie'' (Stg 1:13).}

Lamentaciones dice lo mismo desde las ruinas de Jerusalén, y une los dos nombres de la soberanía:

\hebverso{מִי זֶה אָמַר וַתֶּהִי אֲדֹנָי לֹא צִוָּה \\ מִפִּי עֶלְיוֹן לֹא תֵצֵא הָרָעוֹת וְהַטּוֹב}
\begin{verso}{Lamentaciones 3:37–38}
¿Quién será aquel que diga que sucedió algo que el Señor no mandó? ¿De la boca del Altísimo no sale lo malo y lo bueno?
\end{verso}

\begin{filologia}{Lm 3:37--38}
\heb{אֲדֹנָי} \emph{Adonay} & \heb{אדן} & ``el Señor'' (3:37). \\
\heb{עֶלְיוֹן} \emph{Elyón} & \heb{עלה} & ``el Altísimo'' (3:38). \\
\heb{הָרָעוֹת וְהַטּוֹב} & \heb{רעע} & ``los males y el bien'': las calamidades (la caída de la ciudad) y la prosperidad. \\
\end{filologia}

No es teoría: es fe en medio del dolor. Pocos versos antes, el mismo poeta dice: ``Bueno es \SENOR{} a los que en él esperan'' (3:25) y ``no aflige ni entristece voluntariamente a los hijos de los hombres'' (3:33). Amós lo afirma del juicio sobre la ciudad: ``¿Habrá algún mal en la ciudad, el cual \SENOR{} no haya hecho?'' (Am 3:6). Dios gobierna el mal sin ser autor del pecado; reina también sobre la calamidad.

%===================================================================
\section{La soberanía y el problema del dolor: C. S. Lewis}

\subsection{El libro y su pregunta}
\emph{The Problem of Pain} (1940), traducido al español como \emph{El problema del dolor} (Rialp, 1994), fue uno de los primeros libros apologéticos de C. S. Lewis. No pretende consolar sino pensar: ¿cómo puede un Dios bueno y todopoderoso permitir el sufrimiento? Su respuesta no niega la soberanía de Dios; intenta explicarla, aclarando qué significan ``bueno'', ``todopoderoso'' y ``feliz''. Para este estudio interesan sobre todo cuatro capítulos: ``La omnipotencia divina'' (cap. 2), ``La bondad divina'' (cap. 3) y ``El dolor humano'' (caps. 6--7).

\precision{las citas breves de Lewis provienen de una traducción española disponible (caps. 2--3) o son traducción propia del original inglés (cap. 6). Conviene cotejarlas y fijar las páginas en la edición de Rialp (1994).}

\subsection{Lo que significa “todopoderoso”}
La primera precisión de Lewis es que la omnipotencia no incluye lo intrínsecamente imposible. Un círculo cuadrado no es una ``cosa'' que Dios deje de hacer; es un sinsentido, y ``el disparate sigue siendo disparate aun cuando nos refiramos a Dios'' (Lewis, cap. 2). La Escritura dice algo semejante cuando afirma que Dios ``no puede negarse a sí mismo'' (2 Ti 2:13) y que ``es imposible que Dios mienta'' (He 6:18): su poder es el poder de lo que él es.

La segunda precisión toca el mundo creado. Si Dios quiere criaturas libres que se encuentren y se relacionen, necesita darles un mundo con un orden natural estable: la madera que sostiene una casa también puede golpear una cabeza. Si se intenta excluir la posibilidad del sufrimiento que traen ese orden y la existencia de voluntades libres, dice Lewis, se encuentra que se ha excluido ``la vida misma'' (cap. 2).

\subsection{El dolor que nos causamos}
El punto que más importa a este estudio está en el capítulo 6. Lewis observa que, cuando las almas se vuelven malas, usan la posibilidad del dolor para herirse unas a otras, y estima que esto explica quizás \destacar{cuatro quintas partes} de los sufrimientos humanos. Añade una lista: no fue Dios, sino el ser humano, quien fabricó los potros de tortura, los látigos, las cárceles, la esclavitud, los cañones, las bayonetas y las bombas.

\precision{resumen y paráfrasis del cap. 6 (``Human Pain'') del original. La cifra de las ``cuatro quintas partes'' es una conjetura de Lewis (``quizás''), no una estadística.}

Aquí está la pregunta de este capítulo: ¿cómo es Dios soberano cuando el dolor lo causa el ser humano? Lewis responde que el mal humano es real y culpable, y que Dios no lo suprime a cada paso porque eso anularía la libertad que hace posible el amor. Pero no concluye que Dios haya perdido el gobierno: el dolor sigue en sus manos y sirve a sus fines.

\subsection{El megáfono de Dios}
La frase más conocida del libro está en el mismo capítulo: Dios nos susurra en nuestros placeres, nos habla en nuestra conciencia, pero nos grita en nuestros dolores: el dolor es ``su megáfono para despertar a un mundo sordo'' (Lewis, cap. 6, trad. propia). El dolor no se puede ignorar; rompe la ilusión de que todo está bien y de que nos bastamos a nosotros mismos. Es la experiencia de Nabucodonosor (``hasta que conozcas'', Dn 4:25) y la del salmista:

\hebverso{טֶרֶם אֶעֱנֶה אֲנִי שֹׁגֵג וְעַתָּה אִמְרָתְךָ שָׁמָרְתִּי}
\begin{verso}{Salmo 119:67}
Antes que fuera yo humillado, descarriado andaba; Mas ahora guardo tu palabra.
\end{verso}

\begin{filologia}{Sal 119:67}
\heb{אֶעֱנֶה} \emph{e'ené} & \heb{ענה} & ``fuera afligido, humillado'': la misma raíz de \emph{oní}, ``aflicción''. \\
\heb{שֹׁגֵג} \emph{shoguég} & \heb{שׁגג} & ``andaba errado'', extraviado por ignorancia. \\
\end{filologia}

\subsection{Un Padre, no un abuelo}
En ``La bondad divina'' (cap. 3), Lewis distingue el amor de la mera benevolencia. Queremos un Dios que nos deje hacer lo que nos gusta y nos vea contentos; en sus palabras, ``deseamos no tanto un padre en los cielos, sino más bien un abuelito''. El amor verdadero quiere el bien del amado y no solo su comodidad; por eso puede corregir, y la corrección duele. La carta a los Hebreos lo dice así:

\begin{verso}{Hebreos 12:11}
Es verdad que ninguna disciplina al presente parece ser causa de gozo, sino de tristeza; pero después da fruto apacible de justicia a los que en ella han sido ejercitados.
\end{verso}

\subsection{Lewis y la Escritura: dos acentos}
\begin{center}\small\renewcommand{\arraystretch}{1.2}
\begin{tabular}{@{}L{2.8cm}L{5.0cm}L{6.0cm}@{}}
\toprule
\textbf{Tema} & \textbf{Lewis} & \textbf{Pink y los textos} \\ \midrule
Origen del dolor & la libertad de las criaturas y un orden natural estable & también; pero el mal humano queda dentro del consejo de Dios (Gn 50:20; Hch 4:27--28) \\
Papel de Dios & permite el dolor y lo usa para despertar & ``hace según su voluntad'' (Dn 4:35); ``\SENOR{} dio, y \SENOR{} quitó'' (Job 1:21) \\
Consuelo & el dolor tiene sentido y no es la última palabra & ``todas las cosas les ayudan a bien'' (Ro 8:28); ``eterno peso de gloria'' (2 Co 4:17) \\
\bottomrule
\end{tabular}\par
{\footnotesize\color{gris}\textit{Tabla 7.} Lewis y Pink ante el dolor. \textit{Nota.} Elaboración propia.}
\end{center}

Lewis subraya la libertad humana; Pink, la soberanía que la incluye. La Escritura sostiene ambas: el hombre que hiere es culpable, y Dios sigue reinando sobre ese mismo mal.

\precision{la comparación es propia de esta serie. Lewis escribe como apologista ante un lector escéptico; Pink, como expositor ante creyentes. Sus acentos difieren, pero no se excluyen del todo.}

\subsection{Job: el dolor bajo el trono}
El libro de Job es la versión bíblica del problema. Las pérdidas de Job vienen de sabeos, caldeos, fuego y viento (1:13--19); detrás está el Satán, y por encima, el permiso de Dios (1:12; 2:6). Job no atribuye su dolor a los sabeos ni al Satán:

\hebverso{יְהוָה נָתַן וַיהוָה לָקָח יְהִי שֵׁם יְהוָה מְבֹרָךְ}
\begin{verso}{Job 1:21}
y dijo: Desnudo salí del vientre de mi madre, y desnudo volveré allá. \SENOR{} dio, y \SENOR{} quitó; sea el nombre de \SENOR{} bendito.
\end{verso}
\begin{verso}{Job 2:10}
Y él le dijo: Como suele hablar cualquiera de las mujeres fatuas, has hablado. ¿Qué? ¿Recibiremos de Dios el bien, y el mal no lo recibiremos? En todo esto no pecó Job con sus labios.
\end{verso}

Dios nunca le explica el porqué. Le muestra quién es (Job 38--41), y Job responde: ``Yo conozco que todo lo puedes, y que no hay pensamiento que se esconda de ti'' (42:2).

%===================================================================
\section{Un encuentro en la historia de Israel: José}

\subsection{Trece años de dolor}
José tenía diecisiete años cuando sus hermanos lo odiaron (Gn 37:2, 4) y treinta cuando se presentó ante Faraón (41:46). Entre una edad y otra hay trece años de dolor causado por otros: los hermanos que lo vendieron, la mujer de Potifar que lo calumnió, el copero que lo olvidó (40:23).

\hebverso{וַיִּמְכְּרוּ אֶת־יוֹסֵף לַיִּשְׁמְעֵאלִים בְּעֶשְׂרִים כָּסֶף}
\begin{verso}{Génesis 37:28}
Y cuando pasaban los madianitas mercaderes, sacaron ellos a José de la cisterna, y le trajeron arriba, y le vendieron a los ismaelitas por veinte piezas de plata. Y llevaron a José a Egipto.
\end{verso}

Es el caso que plantea la pregunta de Lewis en forma de historia: ¿dónde está Dios cuando el dolor lo causa el ser humano?

\subsection{“No me enviasteis acá vosotros”}
Cuando José se da a conocer a sus hermanos, interpreta su historia:

\hebverso{וְעַתָּה לֹא־אַתֶּם שְׁלַחְתֶּם אֹתִי הֵנָּה כִּי הָאֱלֹהִים \\ וַיְשִׂימֵנִי לְאָב לְפַרְעֹה וּלְאָדוֹן לְכָל־בֵּיתוֹ וּמֹשֵׁל בְּכָל־אֶרֶץ מִצְרָיִם}
\begin{verso}{Génesis 45:5, 8}
Ahora, pues, no os entristezcáis, ni os pese de haberme vendido acá; porque para preservación de vida me envió Dios delante de vosotros\ldots{} Así, pues, no me enviasteis acá vosotros, sino Dios, que me ha puesto por padre de Faraón y por señor de toda su casa, y por gobernador en toda la tierra de Egipto.
\end{verso}

\begin{filologia}{Gn 45:5--8}
\heb{כִּי־מְכַרְתֶּם אֹתִי} & \heb{מכר} & ``que me vendisteis'': José no niega el pecado de sus hermanos. \\
\heb{שְׁלָחַנִי אֱלֹהִים} & \heb{שׁלח} & ``me envió Dios'': tres veces en tres versículos (45:5, 7, 8). \\
\heb{לְמִחְיָה} \emph{lemijyá} & \heb{חיה} & ``para preservación de vida''. \\
\heb{וּלְאָדוֹן} \emph{uleadón} & \heb{אדן} & ``y por señor'' de la casa de Faraón. \\
\heb{וּמֹשֵׁל} \emph{umoshel} & \heb{משׁל} & ``y gobernador'': participio de \emph{mashal}, ``gobernar''. \\
\end{filologia}

El esclavo vendido termina como \emph{adón} y \emph{moshel}: las mismas palabras de la soberanía, ahora delegadas. La historia de José es una parábola de la soberanía de Dios: el que gobierna es Dios, y lo hace a través de lo que otros hicieron para mal.

\subsection{“Dios lo encaminó a bien”}
Muerto Jacob, los hermanos temen la venganza (50:15). José responde:

\hebverso{הֲתַחַת אֱלֹהִים אָנִי \\ וְאַתֶּם חֲשַׁבְתֶּם עָלַי רָעָה אֱלֹהִים חֲשָׁבָהּ לְטֹבָה \\ לְמַעַן עֲשֹׂה כַּיּוֹם הַזֶּה לְהַחֲיֹת עַם־רָב}
\begin{verso}{Génesis 50:19–20}
Y les respondió José: No temáis; ¿acaso estoy yo en lugar de Dios? Vosotros pensasteis mal contra mí, mas Dios lo encaminó a bien, para hacer lo que vemos hoy, para mantener en vida a mucho pueblo.
\end{verso}

\begin{filologia}{Gn 50:19--20}
\heb{הֲתַחַת אֱלֹהִים אָנִי} & \heb{תחת} & ``¿acaso estoy yo en lugar de Dios?'': José no ocupa el lugar del Juez. \\
\heb{חֲשַׁבְתֶּם עָלַי רָעָה} & \heb{חשׁב} & ``vosotros planeasteis contra mí mal'': \emph{jashav}, ``pensar, tramar, calcular''. \\
\heb{אֱלֹהִים חֲשָׁבָהּ לְטֹבָה} & \heb{חשׁב} & ``Dios lo planeó para bien'': el \emph{mismo verbo}. El sufijo femenino \heb{ָהּ} (``lo'') remite a la misma \emph{ra'á}, el mismo mal. La RVR1960 traduce ``lo encaminó a bien''. \\
\heb{לְהַחֲיֹת עַם־רָב} & \heb{חיה} & ``para mantener con vida a mucho pueblo'': el fin de la soberanía. \\
\textbf{LXX} & --- & \grc{ὑμεῖς ἐβουλεύσασθε κατ’ ἐμοῦ εἰς πονηρά, ὁ δὲ θεὸς ἐβουλεύσατο περὶ ἐμοῦ εἰς ἀγαθά}: ``vosotros deliberasteis contra mí para mal; Dios deliberó acerca de mí para bien''. \\
\end{filologia}

No son dos hechos, uno malo y otro bueno: es \destacar{un solo hecho} con dos intenciones. Los hermanos son culpables, y Dios es soberano. La Septuaginta traduce \emph{jashav} con \grc{βουλεύω}, el verbo de la \grc{βουλή}, el ``designio'' de Dios que Pablo nombra en Ef 1:11. José no ocupa el lugar de Dios para vengarse, porque Dios ya gobernó lo que ellos hicieron.

%===================================================================
\section{La soberanía en el judaísmo}

\precision{esta sección recoge lecturas judías antiguas y medievales. Son interpretaciones rabínicas, citadas en su texto estándar (Misná, Talmud de Babilonia, Maimónides, Qumrán) con traducción propia. En esta sesión no se pudo consultar Sefaria en línea: conviene cotejar cada texto allí.}

\subsection{El yugo del reino de los cielos}
La Misná explica por qué el \emph{Shemá} (Dt 6:4--9) precede, en la oración diaria, a Dt 11:13--21. Rabí Yehoshúa ben Qorjá responde:

\hebverso{כְּדֵי שֶׁיְּקַבֵּל עָלָיו עֹל מַלְכוּת שָׁמַיִם תְּחִלָּה \\ וְאַחַר כָּךְ יְקַבֵּל עָלָיו עֹל מִצְוֹת}

``Para que reciba sobre sí primero el yugo del reino de los cielos, y después reciba sobre sí el yugo de los mandamientos'' (Misná \emph{Berajot} 2:2). Recitar el \emph{Shemá} cada mañana y cada noche es aceptar a Dios como Rey antes de obedecer cualquier mandamiento. A la primera línea del \emph{Shemá} se responde, según una tradición que el Talmud remonta a Jacob (\emph{Pesajim} 56a): \heb{בָּרוּךְ שֵׁם כְּבוֹד מַלְכוּתוֹ לְעוֹלָם וָעֶד}, ``Bendito sea el nombre de la gloria de su reino por siempre jamás''. Y en tiempo de sequía, Rabí Aquiba oró \heb{אָבִינוּ מַלְכֵּנוּ}, ``Padre nuestro, Rey nuestro'' (\emph{Taanit} 25b), fórmula que dio origen a la letanía de los días terribles.

\precision{el Talmud atribuye a Aquiba la fórmula ``Padre nuestro, Rey nuestro''; la letanía larga que se reza hoy es posterior (Elbogen, 1993).}

\subsection{\emph{Maljuyot}: los versos de la realeza}
En la oración adicional de \emph{Rosh Hashaná} hay tres bendiciones centrales (Misná \emph{Rosh Hashaná} 4:5--6): \emph{Maljuyot} (``reinados''), versos que proclaman a Dios Rey; \emph{Zijronot} (``memorias''), versos sobre el Dios que recuerda; y \emph{Shofarot}, versos sobre el cuerno. La Misná pide que no se digan menos de diez versos de realeza, tomados de la Torá, los Profetas y los Escritos. Entre los que se usan están ``\SENOR{} reinará eternamente y para siempre'' (Éx 15:18), ``\SENOR{} reina; se vistió de magnificencia'' (Sal 93:1), ``y el reino será de \SENOR'' (Abd 21) y ``\SENOR{} será rey sobre toda la tierra'' (Zac 14:9). La sección se abre con la oración \emph{Alenu}: ``nos inclinamos, nos postramos y damos gracias ante el Rey de reyes de reyes, el Santo, bendito sea''.

\precision{la lista de versos es ilustrativa; la selección y el orden varían según el rito. Conviene cotejar en un \emph{majzor}.}

\subsection{“Todo está previsto, y la libertad está dada”}
Los sabios sostuvieron las dos verdades de Gn 50:20 sin resolverlas. Rabí Aquiba enseña:

\hebverso{הַכֹּל צָפוּי וְהָרְשׁוּת נְתוּנָה וּבְטוּב הָעוֹלָם נִדּוֹן}

``Todo está previsto, y la libertad está dada; y el mundo es juzgado con bondad'' (Misná \emph{Avot} 3:15). Y Rabí Janiná:

\hebverso{הַכֹּל בִּידֵי שָׁמַיִם חוּץ מִיִּרְאַת שָׁמַיִם}

``Todo está en manos del Cielo, excepto el temor del Cielo'' (\emph{Berajot} 33b), con apoyo en Dt 10:12: ``¿qué pide \SENOR{} tu Dios de ti, sino que temas a \SENOR{} tu Dios\ldots?''. Maimónides lo repetirá como principio: ``la libertad está dada a todo hombre'' (\emph{Hiljot Teshuvá} 5:1). Es el mismo equilibrio de Flp 2:12--13.

\precision{en algunas ediciones el dicho de Aquiba figura como \emph{Avot} 3:19.}

\subsection{El dedo y el decreto}
Rabí Janiná enseña también: \heb{אֵין אָדָם נוֹקֵף אֶצְבָּעוֹ מִלְּמַטָּה אֶלָּא אִם כֵּן מַכְרִיזִין עָלָיו מִלְמַעְלָה}, ``nadie se golpea el dedo aquí abajo si no lo proclaman sobre él desde arriba'' (\emph{Julín} 7b), y cita Sal 37:23: ``Por \SENOR{} son ordenados los pasos del hombre''. Aun el dolor pequeño está bajo el trono.

\subsection{Bendecir por lo malo}
La Misná convierte la fe de Job en liturgia:

\hebverso{חַיָּב אָדָם לְבָרֵךְ עַל הָרָעָה כְּשֵׁם שֶׁהוּא מְבָרֵךְ עַל הַטּוֹבָה}

``Uno debe bendecir por lo malo así como bendice por lo bueno'' (Misná \emph{Berajot} 9:5). Ante una mala noticia se dice \heb{בָּרוּךְ דַּיַּן הָאֱמֶת}, ``Bendito el Juez verdadero'' (\emph{Berajot} 9:2). Es la confesión de Job (``sea el nombre de \SENOR{} bendito'', 1:21) y la de Nabucodonosor (``sus caminos [son] justos'', Dn 4:37).

\subsection{“Todo lo que hace el Misericordioso, lo hace para bien”}
El Talmud transmite una máxima en cadena (Rav Huná en nombre de Rav, en nombre de Rabí Meír, y así se enseñó en nombre de Rabí Aquiba): \heb{כׇּל דְּעָבֵיד רַחֲמָנָא לְטַב עָבֵיד}, ``todo lo que hace el Misericordioso, lo hace para bien'' (\emph{Berajot} 60b). La ilustra con un relato: Aquiba no encuentra posada y duerme en el campo; el viento apaga su lámpara, una comadreja se come su gallo y un león su asno. Esa noche un ejército saquea la aldea; si hubieran visto la luz u oído al gallo o al asno, lo habrían capturado. De Najúm de Guimzó se cuenta que, ante cada desgracia, decía \heb{גַּם זוֹ לְטוֹבָה}, ``también esto es para bien'' (\emph{Taanit} 21a).

\precision{relatos rabínicos resumidos. La relación entre ``para bien'' (\emph{letov}, \emph{letovah}), Gn 50:20 (\emph{letová}) y Ro 8:28 (\grc{εἰς ἀγαθόν}) es lectura propia de esta serie.}

\subsection{Josefo y Qumrán}
En tiempos de Pablo la cuestión dividía a las escuelas judías. Flavio Josefo (\emph{Antigüedades} 13.171--173) describe tres posturas: los fariseos dicen que algunas cosas, no todas, son obra del destino, y que otras dependen de nosotros; los esenios, que el destino gobierna todas las cosas; los saduceos, que no hay destino y todo depende de nosotros. Pablo era fariseo (Flp 3:5).

La \emph{Regla de la Comunidad} de Qumrán, en su instrucción sobre los dos espíritus, dice:

\hebverso{מֵאֵל הַדֵּעוֹת כּוֹל הוֹוֶה וְנִהְיֶה \\ וְלִפְנֵי הֱיוֹתָם הֵכִין כּוֹל מַחֲשַׁבְתָּם}

``Del Dios de los conocimientos procede todo lo que es y lo que será; y antes de que existieran estableció todo su designio'' (1QS 3:15).

\precision{josefo escribe para lectores griegos y usa ``destino'' (\grc{εἱμαρμένη}) donde un judío diría ``providencia''. El texto de 1QS sigue las ediciones de García Martínez (1992) y Vermes (2004); vocalización y traducción propias. Conviene cotejar el manuscrito en la Leon Levy Dead Sea Scrolls Digital Library.}

%===================================================================
\section{La soberanía en Cristo, leída con Pablo}

\subsection{“Según el designio de su voluntad”}
La definición de este capítulo está, casi palabra por palabra, en Efesios:

\begin{verso}{Efesios 1:11}
En él asimismo tuvimos herencia, habiendo sido predestinados conforme al propósito del que hace todas las cosas según el designio de su voluntad.
\end{verso}

\begin{filologia}{Ef 1:11}
\grc{κατὰ πρόθεσιν} & \grc{πρόθεσις} & ``conforme al propósito'': el plan ``puesto delante''. \\
\grc{τοῦ τὰ πάντα ἐνεργοῦντος} & \grc{ἐνεργέω} & ``del que opera todas las cosas'': obra eficazmente (Pink: ``trabajar eficazmente''). \\
\grc{κατὰ τὴν βουλήν} & \grc{βουλή} & ``según el designio'': el sustantivo del verbo de Gn 50:20 LXX. \\
\grc{τοῦ θελήματος αὐτοῦ} & \grc{θέλημα} & ``de su voluntad'': el término de Dn 4:35 en Teodoción. \\
\end{filologia}

Pablo dice lo mismo que Nabucodonosor y que José, pero todo ocurre ``en él'': en Cristo. La soberanía de Dios tiene un centro.

\precision{la autoría paulina de Efesios se discute; la misma idea está en cartas indiscutidas: ``de él, y por él, y para él, son todas las cosas'' (Ro 11:36).}

\subsection{“Todas las cosas les ayudan a bien”}
\begin{verso}{Romanos 8:28}
Y sabemos que a los que aman a Dios, todas las cosas les ayudan a bien, esto es, a los que conforme a su propósito son llamados.
\end{verso}

\begin{filologia}{Ro 8:28}
\grc{πάντα συνεργεῖ} & \grc{συνεργέω} & ``todas las cosas obran juntas'': las cosas no son buenas en sí; Dios las hace cooperar. \\
\grc{εἰς ἀγαθόν} & \grc{ἀγαθός} & ``para bien'': cf. Gn 50:20 LXX, \grc{εἰς ἀγαθά}. \\
\grc{κατὰ πρόθεσιν} & \grc{πρόθεσις} & ``conforme a su propósito'', como en Ef 1:11. \\
\end{filologia}

El contexto es el dolor: ``las aflicciones del tiempo presente'' (8:18), la creación que ``gime a una'' (8:22), nosotros que gemimos (8:23), el Espíritu que intercede ``con gemidos indecibles'' (8:26). El ``bien'' es ser ``hechos conformes a la imagen de su Hijo'' (8:29).

\precision{el eco de Gn 50:20 en Ro 8:28 es posible, no seguro: ambos comparten \grc{εἰς ἀγαθ-} y la idea de un designio de Dios. Es lectura propia.}

\subsection{El alfarero: Romanos 9}
En Romanos 9 Pablo enfrenta la objeción más dura:

\begin{verso}{Romanos 9:19–21}
Pero me dirás: ¿Por qué, pues, inculpa? porque ¿quién ha resistido a su voluntad? Mas antes, oh hombre, ¿quién eres tú, para que alterques con Dios? ¿Dirá el vaso de barro al que lo formó: ¿Por qué me has hecho así? ¿O no tiene potestad el alfarero sobre el barro, para hacer de la misma masa un vaso para honra y otro para deshonra?
\end{verso}

Responde con el lenguaje de Isaías (``¿Dirá el barro al que lo labra: ¿Qué haces?'', Is 45:9; cf. 29:16) y de Daniel (``no hay quien\ldots{} le diga: ¿Qué haces?'', Dn 4:35). Y antes ha citado Éx 33:19: ``Tendré misericordia del que yo tenga misericordia'' (Ro 9:15). La soberanía de Romanos 9 es, ante todo, soberanía \destacar{de misericordia}: ``no depende del que quiere, ni del que corre, sino de Dios que tiene misericordia'' (9:16).

\precision{el contexto de Ro 9--11 es el lugar de Israel y de las naciones en el plan de Dios; su aplicación a la salvación individual se discute entre intérpretes (Moo, 1996; Wright, 2013). Pink lo lee como despliegue de la soberanía absoluta de Dios (2020, p. 34).}

\subsection{“¡Oh profundidad!”}
Los capítulos 9--11 terminan en doxología:

\begin{verso}{Romanos 11:33–36}
¡Oh profundidad de las riquezas de la sabiduría y de la ciencia de Dios! ¡Cuán insondables son sus juicios, e inescrutables sus caminos! Porque ¿quién entendió la mente del Señor? ¿O quién fue su consejero? ¿O quién le dio a él primero, para que le fuese recompensado? Porque de él, y por él, y para él, son todas las cosas. A él sea la gloria por los siglos. Amén.
\end{verso}

Ro 11:34 cita Is 40:13 (``¿Quién enseñó al Espíritu de \SENOR, o le aconsejó enseñándole?''), y 11:35 recoge Job 41:11 [41:3] (``¿Quién me ha dado a mí primero, para que yo restituya?''). Las tres preposiciones (\grc{ἐξ αὐτοῦ καὶ δι’ αὐτοῦ καὶ εἰς αὐτόν}) dicen el origen, el medio y la meta de todas las cosas. La soberanía, en Pablo, no termina en una discusión, sino en adoración.

\subsection{“Jesucristo es el Señor”}
El texto paulino más importante para la soberanía usa Isaías 45, el mismo capítulo de ``yo hago todo esto'':

\hebverso{כִּי־לִי תִּכְרַע כׇּל־בֶּרֶךְ תִּשָּׁבַע כׇּל־לָשׁוֹן}
\begin{verso}{Isaías 45:23}
Por mí mismo hice juramento, de mi boca salió palabra en justicia, y no será revocada: Que a mí se doblará toda rodilla, y jurará toda lengua.
\end{verso}
\begin{verso}{Filipenses 2:9–11}
Por lo cual Dios también le exaltó hasta lo sumo, y le dio un nombre que es sobre todo nombre, para que en el nombre de Jesús se doble toda rodilla de los que están en los cielos, y en la tierra, y debajo de la tierra; y toda lengua confiese que Jesucristo es el Señor, para gloria de Dios Padre.
\end{verso}

\begin{center}\small\renewcommand{\arraystretch}{1.18}
\begin{tabular}{@{}L{2.4cm}L{11.4cm}@{}}
\toprule
\textbf{Texto} & \textbf{Original} \\ \midrule
Is 45:23 (TM) & \heb{כִּי־לִי תִּכְרַע כׇּל־בֶּרֶךְ תִּשָּׁבַע כׇּל־לָשׁוֹן} \\
Is 45:23 (LXX) & \grc{ἐμοὶ κάμψει πᾶν γόνυ καὶ ἐξομολογήσεται πᾶσα γλῶσσα τῷ θεῷ} \\
Flp 2:10--11 & \grc{πᾶν γόνυ κάμψῃ\ldots{} καὶ πᾶσα γλῶσσα ἐξομολογήσηται ὅτι κύριος Ἰησοῦς Χριστός} \\
\bottomrule
\end{tabular}\par
{\footnotesize\color{gris}\textit{Tabla 8.} Is 45:23 en Filipenses 2. \textit{Nota.} BHS, Rahlfs-Hanhart y NA28.}
\end{center}

En Isaías 45 habla el Dios que dice ``no hay más que yo'' (45:6) y ``yo \SENOR{} soy el que hago todo esto'' (45:7). Pablo aplica ese texto de soberanía exclusiva a Jesús. El ``nombre sobre todo nombre'' es \grc{κύριος}, la palabra que la Septuaginta usa por \emph{Adonay} y por el Nombre. El Dueño de Gn 15:2 tiene ahora un nombre humano: Jesús.

\precision{la lectura de \grc{κύριος} como el Nombre divino en Flp 2:9--11 es la mayoritaria (Hays, 1989; Wright, 2013). Pablo usa Is 45:23 también en Ro 14:11.}

\subsection{Reinar hasta el fin}
\begin{verso}{1 Corintios 15:25–28}
Porque preciso es que él reine hasta que haya puesto a todos sus enemigos debajo de sus pies. Y el postrer enemigo que será destruido es la muerte. Porque todas las cosas las sujetó debajo de sus pies\ldots{} para que Dios sea todo en todos.
\end{verso}

Pablo une dos salmos. El v.~25 usa Sal 110:1 (``hasta que ponga a tus enemigos por estrado de tus pies''); el v.~27 cita Sal 8:6 [8:7] (``todo lo pusiste debajo de sus pies''). El señorío que Dios dio al ser humano en la creación, y que este perdió, lo ejerce ahora Cristo. El último enemigo es la muerte. La soberanía de Dios en Cristo termina con el dolor: ``ya no habrá muerte, ni habrá más llanto, ni clamor, ni dolor'' (Ap 21:4).

\subsection{“El bienaventurado y solo Soberano”}
\begin{verso}{1 Timoteo 6:15–16}
la cual a su tiempo mostrará el bienaventurado y solo Soberano, Rey de reyes, y Señor de señores, el único que tiene inmortalidad, que habita en luz inaccesible; a quien ninguno de los hombres ha visto ni puede ver, al cual sea la honra y el imperio sempiterno. Amén.
\end{verso}

\grc{ὁ μακάριος καὶ μόνος δυνάστης} es ``el dichoso y único Potentado''. ``Señor de señores'' viene de Dt 10:17, y el Apocalipsis lo aplica al Cordero (Ap 17:14; 19:16). ``El imperio sempiterno'' (\grc{κράτος αἰώνιον}) recuerda el ``dominio sempiterno'' de Dn 4:34.

\precision{la autoría paulina de 1 Timoteo se discute; la doxología coincide con Ro 11:36 y 16:27, cartas indiscutidas.}

\subsection{La cruz: el caso máximo de Gn 50:20}
Lucas, compañero de Pablo, recoge la oración de la iglesia amenazada:

\begin{verso}{Hechos 4:24, 27–28}
Soberano Señor, tú eres el Dios que hiciste el cielo y la tierra, el mar y todo lo que en ellos hay\ldots{} Porque verdaderamente se unieron en esta ciudad contra tu santo Hijo Jesús, a quien ungiste, Herodes y Poncio Pilato, con los gentiles y el pueblo de Israel, para hacer cuanto tu mano y tu consejo habían antes determinado que sucediera.
\end{verso}

La iglesia llama a Dios \grc{Δέσποτα}, el título de Gn 15:2 en la Septuaginta. Y lee la cruz con la lógica de José: el peor mal humano, el asesinato del Justo, cumple ``tu mano y tu consejo'' (\grc{βουλή}). Pedro lo había dicho en Pentecostés: ``entregado por el determinado consejo y anticipado conocimiento de Dios, prendisteis y matasteis por manos de inicuos'' (Hch 2:23). En la cruz el dolor causado por el hombre alcanzó a Dios mismo, y ahí se ve la soberanía en su forma más alta: el Rey reina desde una cruz y convierte la injusticia en redención. La soberanía no paraliza a la iglesia: le hace pedir ``denuedo'' para hablar (4:29).

\subsection{El aguijón: el dolor de Pablo}
\begin{verso}{2 Corintios 12:7–9}
Y para que la grandeza de las revelaciones no me exaltase desmedidamente, me fue dado un aguijón en mi carne, un mensajero de Satanás que me abofetee, para que no me enaltezca sobremanera; respecto a lo cual tres veces he rogado al Señor, que lo quite de mí. Y me ha dicho: Bástate mi gracia; porque mi poder se perfecciona en la debilidad.
\end{verso}

\begin{filologia}{2 Co 12:7--9}
\grc{ἐδόθη μοι} & \grc{δίδωμι} & ``me fue dado'': pasivo; el que da es Dios, aunque el instrumento sea ``un mensajero de Satanás''. \\
\grc{σκόλοψ τῇ σαρκί} & \grc{σκόλοψ} & ``aguijón, estaca en la carne''. \\
\grc{ἵνα μὴ ὑπεραίρωμαι} & \grc{ὑπεραίρω} & ``para que no me enaltezca'': la lección de Nabucodonosor, aplicada a un apóstol. \\
\grc{ἀρκεῖ σοι ἡ χάρις μου} & \grc{ἀρκέω} & ``te basta mi gracia'': Dios no quita el dolor; da gracia en él. \\
\end{filologia}

Como en Job, el dolor tiene un instrumento hostil y un Señor que lo gobierna. Pablo puede decir entonces: ``esta leve tribulación momentánea produce en nosotros un cada vez más excelente y eterno peso de gloria'' (2 Co 4:17).

\subsection{La soberanía cumplida en Cristo}
Pablo, judío formado en las Escrituras, lee la soberanía de \emph{Adonay} y de \emph{Elyón} cumplida en Cristo. Todo ocurre ``en él'' según el designio de Dios (Ef 1:11); la cruz está dentro de su consejo (Hch 4:28); Jesús recibe el Nombre sobre todo nombre (Flp 2:11); reina hasta vencer la muerte (1 Co 15:25). ``Todas las cosas en él subsisten'' (Col 1:17), y ``en él habita corporalmente toda la plenitud de la Deidad'' (Col 2:9). El trono que Nabucodonosor reconoció tiene hoy un Rey con nombre: Jesús, el Señor.

%===================================================================
\section{Vivir bajo el Dios soberano}

\begin{enumerate}
\item \destacar{Humildad.} Mis logros no son ``la gran Babilonia que yo edifiqué'' (Dn 4:30). ``Él puede humillar a los que andan con soberbia'' (4:37).
\item \destacar{Planes en la mano de Dios.} ``En lugar de lo cual deberíais decir: Si el Señor quiere, viviremos y haremos esto o aquello'' (Stg 4:15). El que planea su camino sabe que \SENOR{} endereza sus pasos (Pr 16:9).
\item \destacar{Oración con valor.} La iglesia de Hechos 4 ora al ``Soberano Señor'' no para huir, sino para hablar con denuedo (Hch 4:24, 29).
\item \destacar{Paz en el dolor causado por otros.} José nombra el mal (``me vendisteis'', Gn 45:5), no ocupa el lugar de Dios (50:19) y mira más lejos (``para mantener en vida a mucho pueblo'', 50:20). Pablo lo traduce en conducta: ``No os venguéis vosotros mismos\ldots{} No seas vencido de lo malo, sino vence con el bien el mal'' (Ro 12:19, 21).
\item \destacar{Responsabilidad.} ``Ocupaos en vuestra salvación con temor y temblor, porque Dios es el que en vosotros produce así el querer como el hacer'' (Flp 2:12--13).
\item \destacar{Adoración.} ``De él, y por él, y para él, son todas las cosas'' (Ro 11:36). Bendecir también en el dolor: ``sea el nombre de \SENOR{} bendito'' (Job 1:21).
\end{enumerate}

Spurgeon tenía razón: ningún atributo consuela tanto en la prueba. Pero el consuelo no consiste en entender cada dolor, sino en saber que el trono no está vacío.

%===================================================================
\section{Síntesis}

``Él hace según su voluntad\ldots{} y no hay quien detenga su mano'' (Dn 4:35). Dios gobierna absolutamente todas las cosas según su voluntad: es \emph{Adonay}, el Dueño, y \emph{Elyón}, el Altísimo. Su soberanía alcanza la naturaleza, las naciones, los corazones y aun el mal que el ser humano causa, sin quitarle la culpa: lo que los hombres planean para mal, Dios lo planea para bien (Gn 50:20). El dolor, que Lewis llamó su megáfono, enseña al soberbio ``que el Altísimo gobierna''. En Cristo, la soberanía tiene rostro: el Crucificado recibe el Nombre sobre todo nombre (Flp 2:9--11) y reina hasta vencer la muerte (1 Co 15:25--26).

\begin{center}\small\renewcommand{\arraystretch}{1.18}
\begin{tabular}{@{}L{2.8cm}L{11.0cm}@{}}
\toprule
\textbf{Mirada} & \textbf{Aporte al atributo} \\ \midrule
Hebrea & \emph{Adonay}, \emph{Elyón}, \emph{shallit}; Daniel 4 y su estribillo; Is 45:7; José y \emph{jashav} (Gn 50:20). \\
Judía & yugo del reino de los cielos; \emph{Maljuyot}; ``todo está previsto, y la libertad está dada''; bendecir por lo malo; ``para bien''; 1QS 3:15. \\
Pink y Lewis & la soberanía como ejercicio de la supremacía; responsabilidad fundada en la soberanía; el dolor, la libertad y el megáfono. \\
Paulina & Ef 1:11; Ro 8:28; Ro 9 y 11; Flp 2:11 (Is 45:23); 1 Co 15:25--28; Hch 4:24--28; el aguijón (2 Co 12:7--9). \\
\bottomrule
\end{tabular}\par
{\footnotesize\color{gris}\textit{Tabla 9.} Aportes a la soberanía de Dios. \textit{Nota.} Elaboración propia.}
\end{center}

%===================================================================
\section*{Glosario del capítulo}
\begin{description}\setlength{\itemsep}{2pt}
\item[\emph{Adón}] (\heb{אָדוֹן}). Señor, amo, dueño (Gn 45:8; Mal 1:6).
\item[\emph{Adonay}] (\heb{אֲדֹנָי}). ``El Señor'', título de Dios; lectura judía del Nombre.
\item[\emph{Boulē}] (\grc{βουλή}). Consejo, designio (Hch 2:23; 4:28; Ef 1:11).
\item[\emph{Despótēs}] (\grc{δεσπότης}). Amo con autoridad plena; ``Soberano'' (Hch 4:24; LXX Gn 15:2).
\item[\emph{Dynástēs}] (\grc{δυνάστης}). Potentado, soberano (1 Ti 6:15).
\item[\emph{Elyón}] (\heb{עֶלְיוֹן}). El Altísimo (Gn 14:19; Dt 32:8; Sal 47:2).
\item[\emph{Illayá}] (\heb{עִלָּיָא}). El Altísimo, en arameo (Dn 4:17, 34).
\item[\emph{Jafets}] (\heb{חָפֵץ}). Querer, complacerse (Sal 115:3; 135:6).
\item[\emph{Jashav}] (\heb{חָשַׁב}). Pensar, planear (Gn 50:20; Pr 16:9).
\item[\emph{Kýrios}] (\grc{κύριος}). Señor; en la LXX, por \emph{Adonay} y el Nombre (Flp 2:11).
\item[\emph{Maljuyot}] (\heb{מַלְכֻיּוֹת}). Versos de la realeza de Dios (Misná \emph{Rosh Hashaná} 4:5--6).
\item[\emph{Mashal}] (\heb{מָשַׁל}). Gobernar, dominar (Sal 103:19; 8:6).
\item[\emph{Ol maljut shamayim}] (\heb{עֹל מַלְכוּת שָׁמַיִם}). Yugo del reino de los cielos (Misná \emph{Berajot} 2:2).
\item[\emph{Próthesis}] (\grc{πρόθεσις}). Propósito (Ro 8:28; Ef 1:11).
\item[\emph{Shallit}] (\heb{שַׁלִּיט}). Que tiene dominio, soberano (Dn 4:17, 25).
\item[\emph{Sholtán}] (\heb{שָׁלְטָן}). Dominio, imperio (Dn 4:34).
\end{description}

